\subsolution{Айнымалы (?) ток}{3.5}

\textbf{Назар аударыңыз:} стандартты шрифтпен автордың бастапқы әдісі сипатталған; {\color{violet}күлгін шрифтпен балама әдіс сипатталған.}

\criterion[type=remark]
{Балама әдіс үшін критерийлер фазалық диаграммаларды \textbf{(ФД)} сипаттау критерийлеріне баламалы. Жалпы балл есептің екі бөлігінің әрқайсысы үшін екі әдістің ішіндегі ең жоғарысына қойылады. Сандық есептеулерде қателіктің көшірілуі жоқ.}

Тізбек \figref{1}-суретте көрсетілгендей, $LC$-бөлігіне $U(t)$ периодты импульс әсер етіп тұрғандай көрінеді. Әрине, бұл схема сыртқы кернеудің идеалдандырылған берілуін ескереді, мұнда $U(t)$ токты өзіне ``алып алмайды'' --- әйтпесе конденсатор заряды бірден $q(t)=CU(t)$ болар еді!

\cpic{0.42}{figures/1b.2.png}[Ұзақтығы $\tau$ және периоды $T$ қысқа импульстер түріндегі кернеу]

$LC$-контурының тербеліс периоды мектеп бағдарламасынан белгілі Томсон формуласына тең:
\begin{eq}[points=0.2, tail={; $\omega=\sqrt{LC}$ қабылданады}]
    T_0 = 2\pi\sqrt{LC}.
\end{eq}

Сандық тұрғыдан алғанда, келесі қатынас өте жоғары дәлдікпен орындалатынына назар аударайық:

\begin{eq}[points=0.1]
    \frac T{T_0} = \frac14
\end{eq}

Бұл бізге есепті әрі қарай шешу үшін қажет болады. $U_0\tau$ ``импульсінің'' әсері $LC$-бөлігіндегі $I$ ток күшінің күрт секірісі ретінде көрінеді. Шынында да, Кирхгофтың екінші ережесі бойынша келесі өрнекті аламыз:

\begin{eq}[points=0.1, tail={; оң жақта $U(t)\Leftrightarrow U_0\Leftrightarrow0$ қабылданады}]
    L\cdot\frac{dI}{dt} + \frac qC = U(t),
\end{eq}

және $\tau$ аралығында интегралдай отырып, ток күшінің шекті болуына байланысты біз $q$ зарядының күрт секірісіне ие бола \textbf{алмайтынымызды} көреміз, сондықтан \eqref{1} өрнегін интегралдау келесідей жеңілдейді:

\begin{eq}[points=0.2]
    \Delta I = \frac{U_0\tau}L.
\end{eq}

Бұл есептегі ең ыңғайлы әдіс --- жүйенің $(q, I)$ фазалық кеңістігіндегі эволюциясын қарастыру. Конденсатордағы және индуктивтіліктегі кернеу тербелістері (күрт секірістерді есептемегенде) $\pi/2$ салыстырмалы фазалық ығысуымен таза синусоидалы болғандықтан, жүйенің ``күй нүктесі'' $(q, I)$ шеңбер бойымен сағат тілімен қозғалады. Бұл жүйе энергиясының өрнегінен көрінеді:
$$
E=\frac{LI^2}{2} + \frac{q^2}{2C},
$$
егер ток осін $I/\sqrt{LC}$ бірліктеріне келтірсек, бұл шеңбер теңдеуі болады. $T$ уақыт аралығы тербеліс периоды $T_0$-ның төрттен бірі болғандықтан, фазалық күй сағат тілімен 90 градусқа бұрылады; $\Delta I$ фазалық нүктенің бір шеңберден екіншісіне ``секіруіне'' мүмкіндік беретін вертикаль вектордың мәніне ие болады. Графиктік шешімі Рис. \figref{2}-де көрсетілген. $t=0$ кезіндегі фазалық күй $(0,0)$ нүктесінде болды; содан кейін нүкте қызыл $\to$ көк $\to$ қызғылт сары $\to$ жасыл траекториялармен жүріп, қайтадан $(0,0)$ күйіне оралды.
\criterion{\textbf{Фазалық диаграммалар әдісі:} (\eqref{6}--\eqref{7} өрнектеріне баламалы келесі 5 пункт)}
\criterion[points=0.1]{\textbf{(ФД)} Фазалық күйдің шеңбер бойымен \textbf{және} сағат тілімен қозғалуы}
\criterion[points=0.1]{\textbf{(ФД)} $I/\sqrt{LC}$ нормализациясы немесе баламасы}
\criterion[points=0.1]{\textbf{(ФД)} Әрбір импульс --- нүктенің орын ауыстыруының тік векторы}
\criterion[points=0.1]{\textbf{(ФД)} \eqref{2} ескерсек, әрбір импульс арасында $\pi/2$-ге бұрылу орын алады}
\criterion[points=0.3]{\textbf{(ФД)} Нүктенің қозғалысы дұрыс салынған (Рис. \figref{2}-дегідей)}
Осыдан $t=3T+\tau$ мезетіндегі энергия

\begin{eq}[points=0.4, prefix={\textbf{(a)} бөліміне жауап}, tail={; негіздемесіз қабылданбайды}]
E=0.
\end{eq}

{\color{violet} \textit{Балама түрде,} заряд пен токтың уақытқа тәуелді эволюциясын ``тікелей'' қарастырайық. Импульстен кейін қорек көзінің кернеуі болмайды және жүйе \eqref{1}-ден алынатын $\omega$ жиілігімен еркін тербелістер жасайды:

\criterion[type=remark]{\textbf{Балама әдіс:} (\textbf{ФД} орнына)}

$$
\omega=\frac1{\sqrt{LC}}.
$$

$t=\tau$ мезетіндегі бастапқы шарттар $q(\tau)=0$, $I(\tau)=\Delta I$, сондықтан $\tau <t<T$ үшін,

\begin{eq}[points=0.2, type=alternative]
    q(t)=\frac{\Delta I}{\omega}\sin(\omega t).
\end{eq}

$T$ уақыттан кейін, бұл тербеліс периодының төрттен бірі болғандықтан, ток күші нөлге айналады, ал заряд максималды болып, $\Delta I/\omega$-ға тең болады. Екінші импульстен кейін ток күші $I(T+\tau) = \Delta I$ болады, демек $T+\tau<t<2T$ аралығындағы $q(t)$ тербелістерінің теңдеуі

\begin{eq}[points=0.5, type=alternative]
    q(t)=\frac{\Delta I}\omega\ab(\sin\omega(t-T)+\cos\omega(t-T))=\frac{\sqrt2\Delta I}\omega\sin\ab(\omega (t-T)+\frac\pi4).
\end{eq}

$t=2T$ мезетіне қарай $q(2T)=q(T+\tau)=\Delta I/\omega$, ал $I(2T)=-I(T+\tau)=-\Delta I$ екенін көреміз, сондықтан үшінші импульстен кейін $I(2T+\tau)=0$. Жүйе $2T+\tau<t<3T$ процесі үшін ток жоқ және заряд максималды күйге ауысты --- демек, ширек периодтан кейін ол токтың модулі максималды $I=-\Delta I$ және заряд жоқ күйге өтеді. Төртінші импульстен кейін біз \eqref{5} туралы қорытынды жасай аламыз.
}
Максималды заряд көк траекторияның сыртқы шеңбердегі көлденең осьпен қиылысу нүктесіне тең. Сыртқы шеңбердің радиусы, көріп отырғанымыздай, вертикаль $\Delta I\sqrt{LC}$ векторының ұзындығынан $\sqrt2$ есе үлкен. Бұл бізді бірден жауапқа әкеледі:

\begin{eq}[points=0.3, tail={; сыртқы шеңбер ішкі шеңберден $\sqrt2$ есе үлкен екендігі немесе екінші импульстен кейінгі энергия $2^{1/4}$ есе көп екендігі қабылданады}]
    q = \sqrt2\cdot\Delta I\sqrt{LC}.
\end{eq}

\criterion[points=-0.2, type=penalty]{$\sqrt{LC}$ көбейткіші ұмытылған}

{\color{violet} Егер \eqref{7} өрнегіне оралсақ, $t=1.5T$ мәнін қою арқылы дәл осындай қорытынды жасауға болады.} \eqref{4} өрнегін ескере отырып, соңғы нәтижені аламыз:

\begin{eq}[points=0.3, numpoints=0.1, prefix={\textbf{(a)} бөліміне жауап}]
    q = U_0\tau\sqrt\frac{2C}{L}=2.09\pu{мКл}.
\end{eq}

\textbf{(b)} пунктінің фазалық диаграммалар әдісімен шешімі де өте қарапайым әрі жылдам. $q_*$ фазалық кеңістіктегі шеңбердің радиусын анықтайтынын көреміз. Импульс берілген кезде жүйенің энергиясы өзгермейтіндіктен, вертикаль $\Delta I$ векторы осы шеңбердің секторлық бұрышы $\pi/2$ болатын хордасы болуы керек. $t=0$ мезетінде нүкте $(q_*,0)$ орнында болғандықтан, секіріс мезетіне жету үшін оған бұрышы $\pi/4$ болатын доғаны жүріп өту қажет. {\color{violet} Бұл сонымен қатар ``уақыт бойынша артқа қарау'' арқылы да көрінеді: $t_*-T$ мезетіндегі жүйенің күйі $t_*+\tau$ мезетіндегі күйді қайталауы тиіс.} Демек,
\criterion{\textbf{Фазалық диаграммалар әдісі:} (келесі 1 пункт \eqref{11}-ге баламалы)}
\criterion[points=0.3]{\textbf{(ФД)} $I/\sqrt{LC}$ векторы $q_*$ шеңберіне 90 градус бұрышпен хорда болып тірелуі тиіс екендігі көрсетілген}

\begin{eq}[points=0.3, prefix={\textbf{(b)} бөліміне жауап}, tail={; сандық мәні болмағаны үшін айыппұл жоқ}]
    t_* = \frac T2=125\pu{мс}.
\end{eq}

Енді $q_*$ туралы. {\color{violet} Қорек көзінің жұмысы $A = U_0\Delta q$ толық $E$ энергиясының сақталуына байланысты нөлге тең екенін байқауға болады. Бұл ток күшінің модулі сақталып, оның таңбасы өзгерген жағдайда ғана мүмкін: $q(t)=q_*\cos\omega t$ өрнегінен бастап, бізде:}

\criterion[type=remark]{\textbf{Балама әдіс:} \textbf{ФД} орнына}

\begin{eq}[points=0.3, type=alternative]
-\omega q_*\sin\omega t_* = -\frac{\Delta I}{2},
\end{eq}

{\color{violet} өйткені секірістен кейін ток күші $\Delta I/2$ болады. $t_*$ мәнін \eqref{11}-ге қойып, $q_*$-ке қатысты шешуге болады,} немесе Рис. \figref{2}-дегідей салу арқылы геометриялық жолмен келесіні табамыз:

\begin{eq}
    [points=0.4, numpoints=0.1, prefix={\textbf{(b)} бөліміне жауап}, tail={; егер $\sqrt{LC}$ көбейткіші ұмытылса, балл қойылмайды}]
    q_* = \frac1{\sqrt2}\cdot\Delta I\sqrt{LC}=U_0\tau\sqrt\frac C{2L}=1.04\pu{мКл}.
\end{eq}

\cpic{1}{figures/1b.3}[\textbf{(a)} (сол жақта) және \textbf{(b)} (оң жақта) пункттеріне арналған фазалық диаграммалар]
